\documentclass[10pt,a4paper]{dinbrief}
\usepackage[utf8]{inputenc}
\usepackage{ngerman}
\usepackage{a4wide}
\usepackage{eurosym}
\usepackage{fancyhdr}
\usepackage{pdfpages}
\address{% 
	Dr. Tobias Griebe\\
	Gartenstra{\ss}e 29\\
	15366 Neuenhagen 
}
\backaddress{Dr. T. Griebe, Gartenstr. 29, 15366 Neuenhagen}
\signature{Dr. Tobias Griebe}
\place{Neuenhagen}
\nowindowrules \nobackaddressrule \nowindowtics \normaladdress
\date{\today}
\pagestyle{empty}
\begin{document} 
\begin{letter}  
{%
	HARTING Stiftung \& Co. KG\\
	Marienwerderstr. 3\\
	D-32339 Espelkamp\\
	Deutschland
}
\subject{\bf Bewerbung Teamleiter IoT System Engineering}
\opening{Sehr geehrte Herr Patrick Lahme,}
auf der Suche einer beruflichen Ver"anderung bin ich mit gro"sem Interesse auf Ihre Stellenausschreibung Teamleiter IoT System Engineering am Standort Berlin gesto"sen. 

In meinen bisherigen T"atigkeiten als Competence Center Leiter bei der adesso AG im sowie zuvor als wissenschaftlicher Mitarbeiter an der Universit"at Duisburg-Essen konnte ich meine Kompetenzen in der F"uhrung von Mitarbeitern sowie bei Konzeption und Entwicklung von digitalen L"osungen im gesamten Spektrum Anforderungsermittlung in fr"uhen Projektphasen bis hin zur Implementierung, Qualit"assicherung und Inbetriebnahme unter Beweis stellen.

Als Leiter der Arbeitsgruppe f"ur mobile Anwendungen und Prozesse am Lehrstuhl f"ur Software Engineering, insb. f"ur mobile Anwendungen der Universit"at Duisburg-Essen hatte ich Gelegenheit meine Kenntnisse in den Bereichen Requirements Engineering, Softwareentwicklung und Qualit"atsmanagement in Forschung und Lehre sowie in zahlreichen Industrieprojekten mit starkem Fokus auf mobile Anwendungen in der Praxis einzusetzen. Die Leitung von IT-Projekten und die Verwendung innovativer Technologien geh"orten dabei ebenso zum Projektalltag wie der Einsatz agiler Arbeitsmethoden und der Aufbau und die fachliche  F"uhrung von Projektteams.

Als Competence Center Leiter bei der adesso AG habe ich den Aufbau F"uhrung eines Teams von Beratern und Entwicklern vorangetrieben. Neben dem Aufbau, der fachlichen und disziplinarischen F"uhrung meines Teams befasse ich mich hier mit Projekten im gesamten Handlungsfeld von Ideation, Identifikation digitalisierbarer Gesch"aftsprozesse, Analyse von Digitalisierungs- und Mobilisierungspotenzial existierender und neuer Fachanforderungen, Requirements Engineering, der Koordination, Durchf"uhrung und Begleitung von IT-Projekten in strategischen und operativen Rollen, Erstellung von Handlungsempfehlungen bei Beschaffung, Implementierung und Inbetriebnahme technischer L"osungen sowie mit der Steuerung interner und externer Entwicklungsdienstleister. 

Bei diesen T"atigkeiten "ubernehme ich im Handlungsfeld Mitarbeiterf"uhrung alle Aktivit"aten zur Entwicklung von Mitarbeitern von Recruiting, Qualifikation, Laufbahnstufenentwicklung bis hin zur Retention, Identifikation von Leistung und Potenzial eigener Mitarbeiter, Ermittlung von Entwicklungsbedarf sowie dem zielgerichteten Einsatz passender Personalentwicklungsinstrumente.

Im Handlungsfeld IT-Projekte bediene ich mich aus einem umfangreichen Methodenbaukasten zur Identifikation und Konkretisierung von Digitalisierungspotenzial (u.~a. Business Model Canvas, Stakeholder-Analyse, Touchpoint-Analyse, Customer Journey Mapping und Customer Safari, Interaction Room) sowie g"angiger Modellierungssprachen und Werkzeuge zur Dokumentation und Analyse von Prozessen (z.~B. BPMN, UML).

Mit Fokus auf die Umsetzung von (IT-)Projekten geh"oren der sichere Umgang mit klassischen und agilen Vorgehensmodellen (insb. Scrum) ebenso zu meinen Kompetenzen wie der Umgang mit State-of-the-Art Werkzeugen zur aktiven und passiven Begleitung von IT-Projekten (z.~B. Jira, Confluence oder MS-TFS). Im Rahmen der von mir verantworteten IT-Projekte konnte ich mir ebenfalls fundierte Kenntnisse in den Bereichen Hardware, Software, Programmiersprachen, Technologien f"ur mobile Softwaresysteme und Cloud-basierte Backend-Systeme aneignen.

Nun bin ich auf der Suche nach einem neuen Engagement in einem innovativen und dynamischen Umfeld und sehe den Herausforderungen des Teamleiter IoT System Engineering mit Spannung entgegen.

Vor diesem Hintergrund bin ich davon "uberzeugt, durch meinen ausgepr"agten technischen und akademischen Hintergrund, meine eigenst"andige und qualit"atsorientierte Arbeitsweise sowie meine F"ahigkeit Mitarbeiter anzuleiten, zu entwickeln und zu motivieren zum Erfolg Ihres Unternehmens beizutragen.

Ich bin voraussichtlich ab Juni 2019 verf"ugbar. %Aufgrund meiner Qualifikation meiner bisherigen Berufsfahrungen liegen meine Gehaltsvorstellungen bei 90.000 Euro/Jahr. 

"Uber eine Einladung zu einem pers"onlichen Gespr"ach freue ich mich sehr.

\closing{Mit freundlichen Gr"u"sen,}
\encl{Lebenslauf, Promotionsurkunde, Diplomurkunde, Zertifikat Licence to Lead F"uhrungskr"afteausbildung, Zertifikat Professional Scrum Product Owner (PSPO I), Zertifikat ISTQB Certified Tester Foundation Level}
\end{letter}
\end{document}